% ============================================================
\chapter{Espaces Sym\'etriques}
\label{ch:symetriques}
% ============================================================

\'Elie Cartan, dans les ann\'ees 1920--1930, entreprend l'une des classifications les plus ambitieuses de la g\'eom\'etrie~: celle des espaces riemanniens dont la courbure est ``la m\^eme en tout point et dans toute direction'' au sens o\`u la sym\'etrie g\'eod\'esique en chaque point est une isom\'etrie globale. Ces \emph{espaces sym\'etriques} englobent les espaces de courbure constante (sph\`eres, espaces hyperboliques), les espaces projectifs, les grassmanniennes, et bien d'autres. Cartan montre que leur classification se ram\`ene \`a celle des alg\`ebres de Lie semi-simples, \'etablissant un pont profond entre g\'eom\'etrie et alg\`ebre. Aujourd'hui, les espaces sym\'etriques apparaissent en physique th\'eorique, en th\'eorie des nombres et en apprentissage automatique (optimisation sur les vari\'et\'es de Stiefel et de Grassmann).

\section{D\'efinition et premiers exemples}

\begin{definition}[Espace sym\'etrique riemannien]
Une vari\'et\'e riemannienne compl\`ete $(M, g)$ est un \emph{espace sym\'etrique
riemannien} si pour tout point $p \in M$, il existe une isom\'etrie involutive
$s_p : M \to M$ telle que~:
\begin{enumerate}
    \item $s_p(p) = p$ (point fixe).
    \item $(ds_p)_p = -\operatorname{Id}_{T_pM}$ (la diff\'erentielle est $-\operatorname{Id}$).
\end{enumerate}
L'isom\'etrie $s_p$ est appel\'ee la \emph{sym\'etrie g\'eod\'esique} en $p$.
\end{definition}

\begin{remarque}
La condition $(ds_p)_p = -\operatorname{Id}$ implique que $s_p$ renverse toutes
les g\'eod\'esiques passant par $p$~: $s_p(\gamma(t)) = \gamma(-t)$.
\end{remarque}

\begin{exemple}[Espaces mod\`eles]
\begin{itemize}
    \item $\R^n$~: $s_p(x) = 2p - x$ (sym\'etrie centrale).
    \item $S^n$~: $s_p$ est la r\'eflexion par rapport au sous-espace passant par $p$
          et le centre.
    \item $\mathbb{H}^n$~: dans le mod\`ele du disque, $s_0$ est l'inversion $x \mapsto -x$.
\end{itemize}
\end{exemple}

\begin{exemple}[Espaces projectifs]
$\R P^n$, $\C P^n$, $\mathbb{H}P^n$ (espace projectif quaternionique) et le plan
de Cayley $\mathbb{O}P^2$ sont des espaces sym\'etriques compacts de rang $1$.
\end{exemple}

\begin{exemple}[Groupes de Lie compacts]
Tout groupe de Lie compact $G$ muni d'une m\'etrique bi-invariante est un espace
sym\'etrique. La sym\'etrie en l'\'el\'ement neutre est $s_e(g) = g^{-1}$.
\end{exemple}

\section{Structure de paire sym\'etrique}

\begin{definition}[Paire sym\'etrique]
Une \emph{paire sym\'etrique riemannienne} est un couple $(G, K)$ o\`u~:
\begin{itemize}
    \item $G$ est un groupe de Lie connexe.
    \item $K$ est un sous-groupe compact de $G$.
    \item Il existe un automorphisme involutif $\sigma : G \to G$ tel que
          $(G^\sigma)_0 \subset K \subset G^\sigma$, o\`u $G^\sigma$ est le
          sous-groupe des points fixes de $\sigma$.
\end{itemize}
\end{definition}

\begin{theoreme}[Correspondance fondamentale]
Il y a une correspondance bijective entre~:
\begin{enumerate}
    \item Les espaces sym\'etriques simplement connexes $(M, g)$.
    \item Les paires sym\'etriques effectives $(G, K, \sigma)$ avec $G$ simplement
          connexe.
\end{enumerate}
L'espace sym\'etrique est $M = G/K$ et la m\'etrique est induite par un produit
scalaire $\operatorname{Ad}(K)$-invariant sur $\mathfrak{p}$.
\end{theoreme}

\section{D\'ecomposition de Cartan}

\begin{definition}[D\'ecomposition de Cartan]
Soit $\sigma : G \to G$ l'involution et $\sigma_* : \mathfrak{g} \to \mathfrak{g}$
sa diff\'erentielle. L'alg\`ebre de Lie se d\'ecompose en espaces propres~:
\[
    \mathfrak{g} = \mathfrak{k} \oplus \mathfrak{p},
\]
o\`u $\mathfrak{k} = \{X \in \mathfrak{g} : \sigma_*(X) = X\}$ (espace propre $+1$)
et $\mathfrak{p} = \{X \in \mathfrak{g} : \sigma_*(X) = -X\}$ (espace propre $-1$).
\end{definition}

\begin{proposition}[Relations de crochets]
Les crochets de Lie satisfont~:
\[
    [\mathfrak{k}, \mathfrak{k}] \subset \mathfrak{k}, \quad
    [\mathfrak{k}, \mathfrak{p}] \subset \mathfrak{p}, \quad
    [\mathfrak{p}, \mathfrak{p}] \subset \mathfrak{k}.
\]
\end{proposition}

\begin{proof}
Si $X, Y \in \mathfrak{k}$, alors $\sigma_*[X,Y] = [\sigma_*X, \sigma_*Y] = [X,Y]$,
donc $[X,Y] \in \mathfrak{k}$. Si $X \in \mathfrak{k}$, $Y \in \mathfrak{p}$,
alors $\sigma_*[X,Y] = [X,-Y] = -[X,Y]$, donc $[X,Y] \in \mathfrak{p}$.
Si $X, Y \in \mathfrak{p}$, alors $\sigma_*[X,Y] = [-X,-Y] = [X,Y]$,
donc $[X,Y] \in \mathfrak{k}$.
\end{proof}

L'espace tangent \`a $M = G/K$ au point base $o = eK$ s'identifie \`a $\mathfrak{p}$,
et la m\'etrique riemannienne est un produit scalaire $\operatorname{Ad}(K)$-invariant
sur $\mathfrak{p}$.

\section{Courbure des espaces sym\'etriques}

\begin{theoreme}[Formule de courbure]
Sur l'espace sym\'etrique $M = G/K$, la connexion de Levi-Civita v\'erifie, pour
$X, Y \in \mathfrak{p} \cong T_oM$~:
\[
    \nabla_X Y = \frac{1}{2}[X, Y]_\mathfrak{p} = 0,
\]
car $[X, Y] \in \mathfrak{k}$ et la projection sur $\mathfrak{p}$ est nulle.
Le tenseur de courbure est~:
\[
    R(X, Y)Z = -[[X, Y], Z], \quad X, Y, Z \in \mathfrak{p}.
\]
\end{theoreme}

\begin{corollaire}
La courbure d'un espace sym\'etrique est enti\`erement d\'etermin\'ee par les crochets
de Lie. En particulier, $\nabla R = 0$ (la courbure est parall\`ele).
\end{corollaire}

\begin{theoreme}[Caract\'erisation par $\nabla R = 0$]
Une vari\'et\'e riemannienne compl\`ete simplement connexe $(M, g)$ est un espace
sym\'etrique si et seulement si $\nabla R = 0$.
\end{theoreme}

\begin{proposition}[Courbure sectionnelle]
Pour $X, Y \in \mathfrak{p}$ orthonorm\'es~:
\[
    K(X, Y) = \inner{[[X,Y], X]}{Y} = \norm{[X, Y]}^2 \geq 0
\]
si la m\'etrique provient de la forme de Killing chang\'ee de signe (type compact).
Pour le type non compact, $K(X,Y) = -\norm{[X,Y]}^2 \leq 0$.
\end{proposition}

\section{Classification des espaces sym\'etriques}

\begin{theoreme}[D\'ecomposition de de Rham]
Tout espace sym\'etrique simplement connexe se d\'ecompose de mani\`ere unique en
produit~:
\[
    M = M_0 \times M_1 \times \cdots \times M_k,
\]
o\`u $M_0 = \R^p$ est le facteur euclidien et les $M_i$ ($i \geq 1$) sont des
espaces sym\'etriques irr\'eductibles.
\end{theoreme}

\begin{definition}[Types d'espaces sym\'etriques]
Les espaces sym\'etriques irr\'eductibles se classent en~:
\begin{itemize}
    \item \textbf{Type compact~:} $K \geq 0$, $G$ compact. Exemple~: $S^n$, $\C P^n$.
    \item \textbf{Type non compact~:} $K \leq 0$, $G$ non compact.
          Exemple~: $\mathbb{H}^n$, $SL(n,\R)/SO(n)$.
    \item \textbf{Dualit\'e~:} \`a tout espace compact correspond un espace non compact
          <<~dual~>>.
\end{itemize}
\end{definition}

\begin{formules}[title=\textbf{Exemples d'espaces sym\'etriques irr\'eductibles}]
\begin{center}
\renewcommand{\arraystretch}{1.3}
\begin{tabular}{lccl}
\hline
\textbf{Espace} & $G/K$ & \textbf{Type} & $\dim$ \\
\hline
Sph\`ere $S^n$ & $SO(n+1)/SO(n)$ & compact & $n$ \\
$\C P^n$ & $SU(n+1)/S(U(1) \times U(n))$ & compact & $2n$ \\
Grassmannienne & $SO(p+q)/S(O(p) \times O(q))$ & compact & $pq$ \\
$\mathbb{H}^n$ & $SO_0(n,1)/SO(n)$ & non compact & $n$ \\
$SL(n,\R)/SO(n)$ & & non compact & $\frac{n(n+1)}{2}-1$ \\
\hline
\end{tabular}
\end{center}
\end{formules}

\section{Rang et sous-espaces plats}

\begin{definition}[Rang]
Le \emph{rang} d'un espace sym\'etrique est la dimension maximale d'un sous-espace
ab\'elien $\mathfrak{a} \subset \mathfrak{p}$ (i.e.\ $[\mathfrak{a}, \mathfrak{a}] = 0$).
G\'eom\'etriquement, c'est la dimension du plus grand sous-espace totalement
g\'eod\'esique plat.
\end{definition}

\begin{exemple}
$S^n$ et $\mathbb{H}^n$ sont de rang $1$. $SL(n,\R)/SO(n)$ est de rang $n-1$.
La Grassmannienne $G_{p,q}$ est de rang $\min(p,q)$.
\end{exemple}

\section{G\'eod\'esiques et structure de groupe}

\begin{proposition}[G\'eod\'esiques]
Les g\'eod\'esiques de $M = G/K$ passant par $o = eK$ sont les courbes
$t \mapsto \exp(tX) \cdot o$ pour $X \in \mathfrak{p}$.
\end{proposition}

\begin{proposition}[Transport parall\`ele]
Le transport parall\`ele le long d'une g\'eod\'esique $\gamma(t) = \exp(tX) \cdot o$
est donn\'e par l'action de $\exp(tX)$ sur $T_oM \cong \mathfrak{p}$.
\end{proposition}

\begin{theoreme}[Holonomie]
Le groupe d'holonomie restreint d'un espace sym\'etrique irr\'eductible est
$\operatorname{Hol}^0(M) = K$ (la composante connexe). La classification de
Berger des groupes d'holonomie des vari\'et\'es irr\'eductibles non sym\'etriques
compl\`ete celle des espaces sym\'etriques.
\end{theoreme}

\section{Espaces sym\'etriques de rang 1}

Les espaces sym\'etriques compacts simplement connexes de rang $1$ sont~:
\begin{enumerate}
    \item $S^n = SO(n+1)/SO(n)$, courbure $K = 1$.
    \item $\C P^n = SU(n+1)/S(U(1) \times U(n))$, courbure $\frac{1}{4} \leq K \leq 1$.
    \item $\mathbb{H}P^n = Sp(n+1)/(Sp(1) \times Sp(n))$, courbure $\frac{1}{4} \leq K \leq 1$.
    \item $\mathbb{O}P^2 = F_4/\operatorname{Spin}(9)$, courbure $\frac{1}{4} \leq K \leq 1$.
\end{enumerate}

\begin{remarque}
Pour $\C P^n$, la courbure sectionnelle varie entre $1/4$ et $1$. Le minimum
$1/4$ est atteint sur les plans totalement r\'eels, le maximum $1$ sur les plans
holomorphes.
\end{remarque}

\section{Exercices}

\begin{exercice}
Montrer que la sph\`ere $S^n$ est un espace sym\'etrique en exhibant explicitement
la sym\'etrie $s_p$ en tout point $p$.
\end{exercice}

\begin{exercice}
Calculer la d\'ecomposition de Cartan $\mathfrak{g} = \mathfrak{k} \oplus \mathfrak{p}$
pour $S^2 = SO(3)/SO(2)$ et v\'erifier les relations de crochets.
\end{exercice}

\begin{exercice}
Montrer que pour le groupe de Lie $G$ muni d'une m\'etrique bi-invariante vu comme
espace sym\'etrique $(G \times G)/\operatorname{diag}(G)$, la courbure sectionnelle
v\'erifie $K(X,Y) = \frac{1}{4}\norm{[X,Y]}^2$.
\end{exercice}

\begin{exercice}
Montrer que $\C P^n$ est un espace sym\'etrique et calculer sa courbure sectionnelle.
V\'erifier que $1/4 \leq K \leq 1$.
\end{exercice}

\begin{exercice}
Montrer que $SL(2,\R)/SO(2) \cong \mathbb{H}^2$ en exhibant une isom\'etrie explicite.
\end{exercice}

\begin{exercice}
D\'eterminer le rang de la Grassmannienne $G_{2,3}(\R) = SO(5)/(SO(2) \times SO(3))$.
\end{exercice}
